\section{Retiming and reverse-complement sectors}

\subsection{Integer retiming and reverse-complement components}
\label{sec:retiming}
The ordinary statements in this section assume $q\ge2$ and $k\ge2$;
the RC statements have the additional hypotheses specified below.
This section works on $G=B(q,k-1)$, with lower vertices and
length-$k$ arcs, including loops. Write $m=\ind_M$ for an ordinary
MDS. A complete move at a residual source $u$ exchanges all
incoming and outgoing arcs; its indicator change is
$\nabla\ind_u$. Inverse moves act at residual sinks.
Let $\CF(M)$ denote the component generated by these moves,
allowing non-RC intermediate MDS.

\begin{lemma}[Classical component interface]\label{lem:move-component}
Two MDS $M,N$ belong to the same ordinary component if and only if
\[
 \ind_N=m+\nabla f\quad\text{for an integer function }f.
\]
Thus $m(Q)=\sum_{e\in Q}m(e)$ is constant on each directed cycle
throughout the component.
\end{lemma}
\begin{proof}
A move adds a unit gradient, and gradients telescope on cycles.
Conversely normalize $\min f=0$. If its maximum $h$ is positive,
choose a source $u$ of the current residual induced on $\{f=h\}$.
An incoming residual arc from a lower level would give
$\ind_N(v\to u)=f(v)-h<0$. Hence $u$ is a source of the entire
residual and a complete move is legal. A new cycle is impossible:
after the exchange no residual arc leaves $u$. For each letter,
$au$ and $ua$ lie in the same PCR, so the one-per-PCR condition
makes the exchange Boolean and preserves cardinality.
A homopolymer loop stays selected and has zero gradient.
Replace $f$ by $f-\ind_u$. Its sum decreases and stays nonnegative,
so iteration reaches $N$.
\end{proof}
The complete-move and cycle-invariant framework is due to
Mar\c cais, DeBlasio and Kingsford
\cite[Sections 4--5]{MarcaisDeBlasioKingsford2024}.
The proof above fixes the gradient sign and the constructive scheduler.

\begin{lemma}[Short separation]\label{lem:short-separation}
Different ordinary components are separated by a simple directed
cycle of length at most $2k-1$ whose selected-arc counts differ.
\end{lemma}
\begin{proof}
Let $d=\ind_N-\ind_M$ and root the graph at $o=a^{k-1}$.
Choose $A_u:o\to u$ by appending $u$ and $B_u:u\to o$ by
appending $a$; both have length at most $k-1$. Put
$p(u)=\sum_{A_u}d$. If some closed walk $A_uB_u$ has nonzero
sum, its length is at most $2k-2$. Otherwise
$\sum_{B_u}d=-p(u)$, and $A_ueB_v$ has sum
$p(u)+d(e)-p(v)$ for $e:u\to v$. If all these sums vanished,
$d=\nabla(-p)$, contradicting \Cref{lem:move-component}.
A nonzero closed walk of length at most $2k-1$ therefore exists.
Decomposing it into simple cycles leaves at least one with
nonzero sum and no greater length.
\end{proof}

\subsubsection{The cycle ratio determines the ordinary optimum}
Define
\[
 R(M)=\max_{Q\text{ simple directed cycle}}\frac{|Q|}{m(Q)}.
\]
Every denominator is positive because $M$ meets every cycle.

\begin{theorem}[Unit retiming]\label{thm:unit-retiming}
For every integer $T\ge1$, the following are equivalent:
\begin{enumerate}
\item Some $N\in\CF(M)$ has $L(N)\le T-1$.
\item There is an integer function $H$ satisfying
\begin{equation}
 H(v)-H(u)\ge1-Tm(e)\qquad(e:u\to v).
 \label{eq:clock}
\end{equation}
\item Every simple directed cycle $Q$ satisfies $|Q|\le Tm(Q)$.
\end{enumerate}
Consequently
\[
 \boxed{\min_{N\in\CF(M)}L(N)=\ceil{R(M)}-1.}
\]
\end{theorem}
\begin{proof}
Given $N$ and $\ind_N=m+\nabla f$, let $h$ be its longest-path
rank, with values between zero and $T-1$. Then $H=h-Tf$
satisfies \eqref{eq:clock}, on both selected and residual arcs.

Conversely divide any integer solution as $H=Ta+h$, with
$0\le h<T$, and set $f=-a$, $m'=m+\nabla f$. The inequality is
\[
 h(v)-h(u)+Tm'(e)\ge1.
\]
Since $h(v)-h(u)\le T-1$, the integer $m'(e)$ is nonnegative.
Each PCR has sum one: the gradient telescopes. The PCRs partition
all arcs, so $m'$ is Boolean. On each zero arc $h$ strictly
increases, giving a DAG of length at most $T-1$ and exactly one
selected arc per PCR. Thus $m'$ is an MDS in the required
component.

Summing \eqref{eq:clock} proves necessity of the cycle inequalities.
For sufficiency give arcs weight $1-Tm(e)$ and add a source with
zero arcs to every vertex. With no positive cycle, a maximum-weight
walk can have its cycles deleted without loss. Maximum walk
weights from the source are finite integers and satisfy
\eqref{eq:clock}. Closed walks decompose into simple cycles.
The smallest feasible integer $T$ is $\ceil{R(M)}$.
\end{proof}
This is a specialization of classical retiming
\cite{LeisersonSaxe1991}. The proof is self-contained; its
MDS-specific step is that nonnegative retiming and one unit per
PCR force Boolean arc values. Bellman--Ford gives
$O(|V||E|)$ integer relaxations per feasibility test; binary search
on $1\le T\le|V|$ finds the optimum. This is polynomial in the
explicit graph, whose size is exponential in $k$, and is distinct
from the direct word-query theorem.

A simple cycle with $r$ selected arcs gives the lower bound
$\ceil{|Q|/r}-1$ throughout a component: its $r$ residual stretches
have total length $|Q|-r$. In particular a one-hit cycle of length
$h+1$ and an attainer of length $h$ certify the component optimum.

\subsubsection{An integer rank with prescribed mate sum}
For the RC results assume $q$ and $k\ge4$ are even, with
fixed-point-free alphabet complement. Then $\theta=\RC$ on lower
vertices is fixed-point-free and reverses arcs.

\begin{lemma}[Bounded autodual rank]\label{lem:autodual-rank}
If a finite DAG has such an involutive antiautomorphism and
longest path at most an integer $c\ge0$, it has an integer rank
$h:V\to\{0,\ldots,c\}$ with
\[
 h(\theta v)=c-h(v),\qquad h(v)-h(u)\ge1\quad(u\to v).
\]
\end{lemma}
\begin{proof}
Let $d(v)$ be the longest path ending at $v$. The half-integral
rank $r(v)=(d(v)+c-d(\theta v))/2$ lies in $[0,c]$, has mate
sum $c$, and increases by at least one on every arc.
Keep its integer values; round half-integers below $c/2$ down
and those above up. For half-integers exactly at $c/2$, choose
opposite rounding in each mate pair. There are no edges between
two such central vertices. For two half-integral endpoints on
the same side the rounding agrees; across the center it
increases the gap, with the worst central choice preserving it.
For one integer endpoint the original gap is at least $3/2$,
and loses at most $1/2$. Thus all final gaps are at least one.
\end{proof}

If $M,N$ are RC and $\ind_N=m+\nabla f$, then
$f(v)+f(\theta v)=C$ is constant, by comparing mate arcs.
Adding an integer constant to $f$ changes $C$ by an even number.
The two parity sectors relative to $M$ are denoted
$\mathcal C_\epsilon(M)$, $\epsilon\in\{0,1\}$, normalizing
$f+f\circ\theta=\epsilon$.

\begin{theorem}[RC retiming in a fixed sector]\label{thm:rc-retiming}
There is $N\in\mathcal C_\epsilon(M)$ with $L(N)\le T-1$ if and
only if an integer solution of \eqref{eq:clock} satisfies
\[
 H(v)+H(\theta v)=c_\epsilon,\qquad
 c_\epsilon=T-1-T\epsilon.
 \tag{RC}
\]
It reconstructs $N$ by $f(v)=-\floor{H(v)/T}$ and
$\ind_N=m+\nabla f$.
\end{theorem}
\begin{proof}
For necessity take the rank from \Cref{lem:autodual-rank} with
$c=T-1$ and set $H=h-Tf$. Conversely write $H=Ta+h$, $0\le h<T$.
Its mate sum is congruent to $-1$ modulo $T$, so the sum of
the two remainders is $T-1$, also when $T=1$. Hence
$a+a\circ\theta=-\epsilon$ and $f+f\circ\theta=\epsilon$.
The ordinary proof gives a Boolean MDS; this constant mate sum
makes its gradient RC invariant.
\end{proof}

Given an ordinary integer solution $d$ of \eqref{eq:clock}, set
\[
 H_*(v)=\frac{d(v)-d(\theta v)+c_\epsilon}{2}.
\]
It is a half-integral solution with the required mate sum. Let
$K_*$ consist of tight arcs
$H_*(v)-H_*(u)=1-Tm(e)$ between half-integral vertices.

\begin{theorem}[Integral feasibility by tight components]
\label{thm:rc-rounding}
The system \eqref{eq:clock} and (RC) has an integer solution
if and only if no strongly connected component of $K_*$
contains both $v$ and $\theta v$.
\end{theorem}
\begin{proof}
If tight paths go both ways between $v$ and $\theta v$, their
weights force every feasible solution to have the same difference
$H(\theta v)-H(v)$ as $H_*$. The mate sum then forces
$2H(v)=2H_*(v)$, an odd integer, a contradiction.

Otherwise the strongly connected components of $K_*$ come in
distinct mate pairs. Number the condensation topologically by $i$.
Round a half-integer up exactly when
$i([v])>i([\theta v])$, otherwise down, leaving integers fixed.
Mate roundings are opposite. A tight edge cannot round its tail
up and head down: that would require
\[
 i([\theta u])<i([u])\le i([v])<i([\theta v])
                    \le i([\theta u]).
\]
A nontight edge with two half-integral endpoints has slack at
least one and loses at most one. With one half-integral endpoint
it has slack at least $1/2$ and loses at most $1/2$.
Thus all inequalities and mate sums survive.
\end{proof}
The tight-component integrality mechanism has a classical UTVPI
antecedent in Lahiri--Musuvathi \cite[Section 4.2]{LahiriMusuvathi2005}.
Here the prescribed mate sum and PCR reconstruction implement it
for Boolean minimum decycling sets.

\begin{corollary}[At most one unit of parity cost]\label{cor:parity-cost}
Let $T_0=\ceil{R(M)}$. Both RC sectors are nonempty and
\[
 \min_{N\in\mathcal C_\epsilon(M)}L(N)\in\{T_0-1,T_0\}.
\]
If $R(M)$ is nonintegral, both optima equal $T_0-1$.
\end{corollary}
\begin{proof}
For $T>R(M)$ every cycle has strictly negative clock weight,
so the tight graph is acyclic. Its components are singletons,
none containing a mate pair. Apply \Cref{thm:rc-rounding} at
$T=T_0$ when noncritical, otherwise at $T=T_0+1$.
The ordinary optimum supplies the lower bound.
\end{proof}
\Cref{prop:sharp-parity} supplies an example attaining the extra
unit in one prescribed sector.

\subsubsection{Exactly two paired-move components}
A paired RC step is a complete move at $u$ and its inverse at
$\theta u$, with disjoint incident supports. RC is required
after the compound step, not after its first half.

\begin{theorem}[Exactly two RC sectors]\label{thm:two-rc-sectors}
Every ordinary component containing an RC MDS has exactly two
components under paired RC steps. They are the two parity sectors.
In one sector, if $\ind_N-\ind_M=\nabla g$ and
$g\circ\theta=-g$, the paired-step distance is
\[
 d_{\rm pair}(M,N)=\tfrac12\sum_v|g(v)|.
\]
Between different sectors the membership of every RC-fixed
length-$k$ word is reversed.
\end{theorem}
\begin{proof}
Both sectors exist by \Cref{cor:parity-cost}. In a single sector
the difference potential has a unique integer antisymmetric
centering $g$. If its maximum is $h>0$, choose a residual source
$u$ on that level as in \Cref{lem:move-component}. Its mate is a
sink at level $-h$. An edge between them would have gradient
$\pm2h$, impossible as the difference of two Boolean indicators.
Their stars are disjoint and the paired move is legal. It replaces
$g$ by $g-\ind_u+\ind_{\theta u}$, reducing its norm by two.
Iteration gives the asserted length; every paired step changes
the norm towards the target by at most two, proving optimality.
Paired steps preserve parity. For a fixed arc $e:u\to\theta u$,
\[
 \ind_N(e)-m(e)=2f(u)-\epsilon.
\]
It is zero for $\epsilon=0$ and $\pm1$ for $\epsilon=1$.
Such words exist at even order, so the sectors are distinct.
\end{proof}
Within the weight-selector family a fixed periodic word
$(a\bar a)^{k/2}$ always has the same frozen zero-PCR membership.
Thus two terminals in the same ordinary component lie in the
same paired sector. The second sector need not be realized by weights.

\begin{proposition}[Components encountered by the DNA family]
\label{prop:family-components}
At $k=4,6,8$ the fixed labelled DNA family $\mathcal W$ with fixed $A<C<G<T$ ties encounters respectively
$1,3,21$ ordinary components, and the same numbers of paired
components. At $k=8$ its component diagram consists of the eleven
intervals and ten distinct points cut out by
\[
\begin{gathered}
 (1+2\sqrt2)/3,\ \sqrt2,\ (5+4\sqrt2)/7,\ (3+6\sqrt2)/7,\
 2\sqrt2-1,\\
 1+\sqrt2,\ 3,\ 1+2\sqrt2,\ 3+2\sqrt2,\ 5+4\sqrt2 .
\end{gathered}
\]
Exactly ten fixed homogeneous linear sign forms suffice and are
necessary to identify this component.
\end{proposition}
\begin{proof}[Exact certificates]
The atlas covers every real ratio $\rho>1$ in $\mathcal W$. Cut-potential normal forms
classify its 199 strata; literal paired moves join each to its
assigned representative. The replay checks 3857 compound steps.
Every pair of distinct classes has a separating cycle: three at
$k=6$ and 210 at $k=8$, all of length at most $2k-1$.
These are positive and negative certificates, not an assumption
that a search found all paths. The component index is
$2\#\{r<\rho\}+\ind_{\rho\text{ is a listed root}}$.
A fixed linear form can have only one normalized root;
each listed wall is indispensable and one lambda form per root
suffices. The periodic anchor proves that the same classification
holds for paired moves inside the family.
\end{proof}
This classifies components encountered by the selector, not all
components in the graph. As a small independent boundary check,
binary order four has six RC MDS in three ordinary components;
all six are isolated under paired moves. Parity cannot be discarded.
